\begin{lemma}[The top homogeneous component of a nonzero polynomial]
Let $\sigma$ be any index type, let $R$ be a commutative semiring and let
$P\in R[X_a:a\in\sigma]$ be nonzero. Then
\[
  H_{\deg P}(P)\ne0,
\]
where $\deg P$ is the total degree of $P$ and $H_d$ is the $d$-th homogeneous
component.
\end{lemma}
\begin{proof}
Since $P\ne0$ its support is nonempty. The total degree of $P$ is by definition
the largest of the entry sums $\sum_a d_a$ over the exponent tuples
$d\in\operatorname{supp}P$; as $\operatorname{supp}P$ is a nonempty finite set,
that largest value is attained, so we may fix $d\in\operatorname{supp}P$ with
\begin{equation}\label{eq:tcnz-attained}
  \sum_a d_a=\deg P .
\end{equation}

Now compute the $d$-th coefficient of the top homogeneous component. By the
definition of the homogeneous component, $H_{d_0}(P)$ keeps exactly the
monomials of $P$ whose exponent tuple has entry sum $d_0$; coefficientwise,
\[
  \operatorname{coeff}_d\bigl(H_{d_0}(P)\bigr)
  =\begin{cases}\operatorname{coeff}_dP,&\sum_a d_a=d_0,\\ 0,&\text{otherwise.}\end{cases}
\]
Taking $d_0=\deg P$ and using \eqref{eq:tcnz-attained}, the first case applies,
so
\begin{equation}\label{eq:tcnz-coeff}
  \operatorname{coeff}_d\bigl(H_{\deg P}(P)\bigr)=\operatorname{coeff}_dP .
\end{equation}

Suppose, for contradiction, that $H_{\deg P}(P)=0$. Then the left-hand side of
\eqref{eq:tcnz-coeff} is the $d$-th coefficient of the zero polynomial, i.e.
$0$, so $\operatorname{coeff}_dP=0$. But $d\in\operatorname{supp}P$ means
precisely that $\operatorname{coeff}_dP\ne0$ --- a contradiction. Hence
$H_{\deg P}(P)\ne0$.
\end{proof}
